% Generated draft: requires mathematical and source review.
\documentclass[11pt]{article}
\usepackage[margin=25mm]{geometry}
\usepackage{fontspec}
\setmainfont{Latin Modern Roman}
\usepackage{amsmath,amssymb,mathtools,unicode-math}
\setmathfont{Latin Modern Math}
\usepackage{enumitem,needspace,xurl,hyperref}
\setlength{\parindent}{0pt}
\setlength{\parskip}{5pt}
\setlength{\emergencystretch}{3em}
\newcommand{\N}{\mathbb{N}}
\newcommand{\Z}{\mathbb{Z}}
\newcommand{\Q}{\mathbb{Q}}
\newcommand{\R}{\mathbb{R}}
\newcommand{\C}{\mathbb{C}}
\newcommand{\F}{\mathbb{F}}
\newcommand{\A}{\mathbb{A}}
\newcommand{\PP}{\mathbb{P}}
\newcommand{\E}{\mathbb{E}}
\newcommand{\eps}{\varepsilon}
\newcommand{\dd}{\,\mathrm{d}}
\begin{document}
\section*{UnsolvedMath ID: 700741}
\textbf{Source problem code:} catalogue entry
\subsection{Definitions and mathematical statement}
Let \(G\) be a group and let \(\mathbb ZG\) denote its integral group ring. A \(\mathbb ZG\)-module \(M\) is Artinian if every descending chain of \(\mathbb ZG\)-submodules stabilizes. Countability refers to the underlying set of \(M\), and includes finite modules.

Define \(G^{(0)}=G\) and \(G^{(i+1)}=[G^{(i)},G^{(i)}]\). The group \(G\) is soluble if \(G^{(d)}=1\) for some integer \(d\geq 0\). Here finite rank is interpreted as finite Pr\"ufer rank: there exists an integer \(r\geq 1\) such that every finitely generated subgroup of \(G\) can be generated by at most \(r\) elements. The supplied source does not explicitly define its rank convention, so this interpretation requires review.

\textbf{Question (L. A. Kurdachenko, 1990; Kourovka Notebook 11.49).} Is it true that, for every soluble group \(G\) of finite rank and every Artinian \(\mathbb ZG\)-module \(M\),
\[
|M|\leq\aleph_0?
\]
The source specifically includes the case in which \(G\) is a soluble minimax group. In the usual soluble-group terminology, a minimax group admits a finite subnormal series whose factors are finite cyclic groups, infinite cyclic groups, or quasicyclic groups \(C_{p^\infty}\), where the prime \(p\) may vary between factors and \(C_{p^\infty}\) is the group of all complex roots of unity of \(p\)-power order.

\textbf{Supplied motivation and qualifications.} The source reports that B. Hartley constructed an uncountable Artinian \(\mathbb ZG\)-module for a metabelian group \(G\) satisfying the minimum condition for normal subgroups. Here metabelian means \(G^{(2)}=1\), and the minimum condition for normal subgroups, denoted \(\mathrm{Min}\text{-}n\), means that every descending chain of normal subgroups stabilizes. The source further reports that this yields an uncountable soluble group of derived length exactly \(3\) satisfying \(\mathrm{Min}\text{-}n\). These are supplied background assertions, not additional hypotheses of the question. The source also connects the question with the study of classes of soluble groups satisfying a weak minimum condition for normal subgroups, without defining that weaker condition.
\subsection{Short English statement}
Must every Artinian module over the integral group ring of a soluble group of finite rank be countable? The question includes soluble minimax groups. Finite rank is interpreted here as finite Pr\"ufer rank, subject to review. The supplied registry labels the problem open; that label is not an independently verified literature assessment.
\subsection{Sources}
\begin{itemize}
\item Supplied local UnsolvedMath registry record: permanent ID \(700741\), exact source code \texttt{catalogue entry}, title \emph{Kourovka Notebook Problem 11.49}. The record attributes the problem to L. A. Kurdachenko, proposed in 1990, and labels it open. Registry creation and update date: 2026-10-16; upstream revision \texttt{d371ee917ac9647fbb06cd6fc529deec6f3dda82}.
\item \emph{The Kourovka Notebook (Unsolved Problems in Group Theory)}, original problem 11.49, 11th issue (1990); supplied reference to the 21st issue edition (2026), arXiv:1401.0300v48. The supplied collection description names E. I. Khukhro and V. D. Mazurov as editors.
\item B. Hartley, \emph{Proceedings of the London Mathematical Society}, volume 35, no. 1 (1977), pp. 55--75. This is the reference supplied for the background construction; the record supplies no article title or URL.
\item These bibliographic references and mathematical background assertions are reproduced from the supplied record and have not been independently checked. No separate literature-verification date is supplied.
\end{itemize}
\end{document}
